\subsection{The first contrast on the original episode window}

\propositionheading{A first-block null under slower source learning}{prop:phil-early-block-null}
In the model of \cref{conj:phil-contrast-formation}, let
\(0<\mathrm T\le10^{-16}\), and let the initial law be
\(\mathcal N(\mathsf m_0,I_{12})\) times uniform assignments.
The learned and frozen arms have the same initial mean and differ only in
the source-learning mobility \(\epsilon_{\mathrm L}\).
Set
\[
 H_{\mathrm T}=2\mathrm T^{-1/8}.
\]
Keep the contrast event of \cref{prop:phil-seeded-contrast-formation}
with its regime parameter equal to \(\mathrm T^{1/8}\): an eligible
proper connected block has an internal belief relative entropy exceeding
\(1/4\) initially, and for a duration-one interval inside
\([0,H_{\mathrm T}]\) all its internal belief and model relative
entropies are at most \(\sqrt{\mathrm T}\), while its boundary belief
relative entropies exceed \(1/2\).
For each block take the infimum of its qualifying completion times,
with value infinity when there is no such interval. Let \(K_{\mathrm T}\)
be the block with least such infimum, using one fixed tie rule in both
arms, and give it a cemetery value when no block qualifies. Write
\(E_{\mathrm T}\) for the event that its value is not the cemetery value.
The infimum need not itself be a qualifying completion time.

Define the positive, temperature-independent constant
\[
 \mathsf c(\mathsf m_0)
 =\frac{e^{-1/4}\varrho_\circ^{12}}{737280}
   \exp\left[-\frac12
       (|\mathbf x^\circ-\mathsf m_0|+\varrho_\circ)^2\right],
 \qquad
 \mathsf d_{\mathrm T}
 =\min\{1,8\epsilon_{\mathrm L}\mathrm T^{-1/4}\}.
\]
Then
\[
 \Pr_0(E_{\mathrm T})\ge\mathsf c(\mathsf m_0),
 \qquad
 \Pr_{\epsilon_{\mathrm L}}(E_{\mathrm T})
       \ge\mathsf c(\mathsf m_0)-\mathsf d_{\mathrm T}.
\]
If \(\mathsf d_{\mathrm T}<\mathsf c(\mathsf m_0)\), the two
conditional block laws satisfy
\[
 \operatorname{TV}\left(
 \operatorname{Law}_{\epsilon_{\mathrm L}}(K_{\mathrm T}\mid E_{\mathrm T}),
 \operatorname{Law}_0(K_{\mathrm T}\mid E_{\mathrm T})\right)
 \le\frac{\mathsf d_{\mathrm T}}{\mathsf c(\mathsf m_0)}.
\]
In particular their difference vanishes for every fixed Gaussian mean
whenever
\(\mathrm T\downarrow0\) and
\(\epsilon_{\mathrm L}\mathrm T^{-1/4}\to0\).
This includes the caps
\(\epsilon_{\mathrm L}\le\mathrm T_n^a\), \(a>7/2\), in
\cref{cor:phil-joint-selection}.
\status{ESTABLISHED}

\noindent\emph{Proof.}
First prove a uniform comparison of the complete native path laws on
any finite horizon \(H\). The learned rows are deterministic and satisfy
\[
 |p_i(t)-1/2|\le\epsilon_{\mathrm L}t/2.
\]
At a common configuration, write \(A,B>0\) for a receiver's likelihood
factors for its two sources. Its posterior positive-source probabilities
in the two arms are
\[
 b(p)=\frac{pA}{pA+(1-p)B},
 \qquad b(1/2)=\frac A{A+B}.
\]
Direct subtraction gives
\[
 |b(p)-b(1/2)|
 =\frac{AB|2p-1|}{(pA+(1-p)B)(A+B)}
 \le|2p-1|.
\]
Indeed the denominator is
\(pA^2+AB+(1-p)B^2\ge AB\). This bound is independent of
temperature and the likelihood ratio, without an interior-row restriction.

Use identical initial samples, a common Brownian motion, and the same
rate-one proposal clock in the two arms. Until the first different redraw,
their configurations and assignments coincide: a held-sheet drift and its
noise do not depend on the prior rows. At each ring, couple each receiver's
two Bernoulli laws maximally. The union bound and the row estimate bound
the probability of any different receiver choice at time \(t\) by
\[
 \sum_i|b_i(p_i(t))-b_i(1/2)|
 \le4\epsilon_{\mathrm L}t.
\]
The compensator of the rate-one clock, stopped at the first mismatch,
bounds its probability before \(H\) by
\(\int_0^H4\epsilon_{\mathrm L}t dt=2\epsilon_{\mathrm L}H^2\).
Consequently,
\begin{equation}
 \operatorname{TV}\left(
 \operatorname{Law}_{\epsilon_{\mathrm L}}(\mathbf X_{[0,H]},S_{[0,H]}),
 \operatorname{Law}_0(\mathbf X_{[0,H]},S_{[0,H]})\right)
 \le\min\{1,2\epsilon_{\mathrm L}H^2\}.
\label{eq:phil-early-path-coupling}
\end{equation}
The bound follows because the complete paths are identical unless the
coupling has had a mismatch. Taking \(H=H_{\mathrm T}\) gives
\(\mathsf d_{\mathrm T}\).

The first-block mark is a measurable function of these paths. To see this
directly, for a fixed block let \(A_I(t)\) be the maximum of its internal
belief and model relative entropies on \([t-1,t]\), and let \(B_I(t)\)
be the minimum boundary belief relative entropy there. Configuration
paths are continuous, so these two window functions are continuous.
Replace the strict condition \(B_I(t)>1/2\) by the countable union of
conditions \(B_I(t)\ge1/2+1/k\), \(k\ge1\). For each \(k\) the
qualifying end times form a closed subset of the compact interval
\([1,H_{\mathrm T}]\). Its first time is measurable, as may be checked
by minimizing
\(\max(A_I(t)-\sqrt{\mathrm T},1/2+1/k-B_I(t))\) on each compact
initial subinterval. The infimum over \(k\), the initial eligibility
condition, and the finite comparison over blocks preserve measurability.
This construction implements the strict boundary convention without
asserting that its infimum is attained.

A measurable map of paths does not increase total variation. Let
\(a_I,b_I\) be the probabilities of the non-cemetery marks in the learned
and frozen arms, \(p=\sum_Ia_I\), \(q=\sum_Ib_I\), and \(d\) the
total-variation distance of the two cemetery-augmented mark laws. Then
\[
 2d=\sum_I|a_I-b_I|+|p-q|,
 \qquad d\le\mathsf d_{\mathrm T}.
\]
If \(p\ge q>0\), the triangle inequality gives
\[
 \frac12\sum_I\left|\frac{a_I}{p}-\frac{b_I}{q}\right|
 \le\frac{\sum_I|a_I-b_I|+(p-q)}{2p}
 =\frac d p.
\]
Interchanging the arms proves the bound \(d/\max(p,q)\).
The frozen matching-seed sub-event in
\cref{prop:phil-seeded-contrast-formation} belongs to \(E_{\mathrm T}\)
and has probability at least \(\mathsf c(\mathsf m_0)\).
Thus \(q\ge\mathsf c(\mathsf m_0)\) and
\(p\ge q-d\ge\mathsf c(\mathsf m_0)-\mathsf d_{\mathrm T}\).
The claimed conditional bound follows, and its limiting conclusion is
immediate. \(\square\)

The proposition keeps the original contrast event and uses its seeded
sub-event only as a probability lower bound. That sub-event already has
the restoring interaction and stationary comparison of
\cref{prop:rg-formation-endogenous-cycle}; the proof does not certify
every history in the broader contrast event under the coherence-episode
rule. Any common measurable admissibility restriction retaining the
seeded sub-event has the same comparison.

The result does not cover the original coupled choice
\(\epsilon_{\mathrm L}=\mathrm T^{1/8}\), for which its bound is
vacuous. It also does not compare different initial means or establish a
conditional parent error. On the slower schedules, late assignment
selection can coexist with a vanishing difference between the conditional
laws of the first contrast on \([0,H_{\mathrm T}]\). These are
different observables on different horizons.


